\documentclass[11pt]{article}
\usepackage{amsmath,amssymb,amsthm}
\newtheorem{theorem}{Theorem}
\title{Signed Residues Repair the Underloaded Thirty-Two-Pole Stress Word}
\author{}
\date{November 24, 2026}
\begin{document}
\maketitle

\paragraph{Status.}
This is an exact finite-dimensional computational certificate, with a
complete rational-function continuation argument. It is not Lean
formalized. It repairs the particular positive-residue obstruction in
\texttt{butterfly-underloaded-h32-obstruction.tex}; it does not prove a
universal decoder, a density estimate, or either global asymptotic target.
No original model definition is changed.

Use the same $H=32$ butterfly, pole grid, retained queries, and routing
as in that obstruction. Explicitly,
\[
 p_y(\epsilon)=\sum_{j=0}^4 y_j\epsilon^{2^j},
 \qquad t_{k,s}=1-p_s(\epsilon).
\]
At switch $e=(k,s,g)$ retain the query exactly when $s<2^k-1$.
There are 49 retained queries. Set
\[
 \theta_{k,s,g}=\left\lfloor s/2^{k-1}\right\rfloor\bmod2
\]
on retained queries, and set every omitted routing bit to zero. Write
$\pi$ for this butterfly permutation, and
$u_{k,g}(h)=1,-1,0$ on the lower half, upper half, and complement
of its input block respectively.

\begin{theorem}
For every sufficiently small $\epsilon>0$, there are real residues
$c_y$ such that every retained query satisfies
\[
 (1-2\theta_{k,s,g})
 \sum_{y=0}^{31}\frac{u_{k,g}(\pi^{-1}(y))c_y}
 {1-p_s(\epsilon)-p_y(\epsilon)}>0.
 \tag{1}
\]
The routing and query set are unchanged. No block biases are needed.
At $\epsilon=1/100$, the saved witness has all 49 signed scores at
least one and has ten negative residue coordinates.
\end{theorem}
\begin{proof}[Certificate and continuation proof]
Let $A(\epsilon)$ be the 49-row signed query matrix in (1).
The companion JSON specifies a square system of 30 query equalities
and the two coordinate equalities $c_0=c_8=0$. The query equalities
have right-hand side one. Its coefficient matrix $B(\epsilon)$ has
entries rational in $\epsilon$.

The exact checker first solves this system over $\mathbb Q$ at
$\epsilon=1/100$, and independently verifies every signed query score.
For continuation it performs guarded elimination on formal series,
with precision 1536 and the right-hand side multiplied by
$\epsilon^{512}$. Every nonzero pivot has a certified leading term;
the pivot orders are
\[
 0,0,0,0,0,0,0,0,0,0,2,2,2,6,6,6,
 8,8,10,16,16,16,18,22,24,41,44,50,53,64,68,72.
\]
Thus $\det B$ is not the zero rational function. The unique rational
solution $c(\epsilon)=B(\epsilon)^{-1}b$ is defined on a sufficiently
small punctured interval. The checker obtains all 49 leading terms
of $A(\epsilon)c(\epsilon)$, with positive coefficients and orders
between $-68$ and $0$. The guards ensure that zero series truncated
before a possible pivot are never used as exact zeros; the final
residual precision is beyond the shifted constant term.

A nonzero rational function with positive Laurent leading coefficient
is positive on a sufficiently small positive interval. Intersect the
finitely many intervals for these 49 scores and the nonzero determinant.
This proves (1).
\end{proof}

\paragraph{Interpretation.}
The earlier exact obstruction included the nonnegativity constraints
on residues $c_2$ and $c_4$ in its positive dependence. It therefore
never excluded this signed solution. Original admissible heads allow
signed numerators. Any attempted universal positive-pool invariant is
strictly stronger than what the original model requires, and its
failure must not be treated as a failure of the signed underloaded
construction.
\end{document}
